\section{Tastenkombinationen}
\label{sec:ref::hotkeys}

Tastaturen haben ein paar besondere Tasten, die es einem erlauben, zum Beispiel Text zu markieren, kopieren und einzufügen. Auch lassen sich dadurch einige spezielle Zeichen, die oft bei der Programmierung verwendet werden, die auf einer deutschen Tastatur keine eigenen Tasten belegen, eingeben. Die Tasten, die ihr dafür benötigt, sind:
\begin{description}
    \item[Shift]
        Die sogenannte ``Umschalttaste''.
        \protect\footnote{%
            Sucht mal auf \url{https://youtube.com} nach «Umschalttaste mechanische Schreibmaschine», wenn ihr wissen wollt, woher dieser Begriff kommt.}%
        Sie schaltet zwischen Klein- und Großschreibung, sowie zwischen den verschiedenen abgebildeten Symbolen um.
    \item[Feststelltaste]
        Im Englischen ``Caps lock'', sie befindet oberhalb der ``Shift''-Taste, sie ist dafür gedacht, dauerhaft zwisch Groß- und Kleinbuchstaben umzuschalten.
        % \footnote{%
        %     Ihr solltet im Hinterkopf behalten, dass diese Taste selten benutzt wird und auch umgedeutet werden kann. Der Autor verwendet diese Taste auf einer Englischen Tastatur, um komplexe Zeichen einzugeben. Das wird als \href{https://de.wikipedia.org/wiki/Compose-Taste}{``Compose-Taste''} bezeichnet, z.B. ergibt bei ihm die Eingabe: CapsLock+"+a das Zeichen ``ä'', oder CapsLock+s+s das Zeichen ``ß''.
        % }
    \item[Strg]
        Die Steuerungstaste (im Englischen ``Control'' bzw. abgekürzt ``Ctrl'').
    \item[Alt]
        ``Alt'' steht für ``Alternative''. Sie wird häufig auch in Verbindung mit ``Shift'' und ``Strg'' verwendet. Allein wird sie z.B.\ verwendet, um zwischen Fenstern mit ``Alt+Tab'' (``Tab'' siehe weiter unten) hin- und herzuschalten.
    \item[AltGr]
        ``AltGr'' steht für ``Alternative Grafik'', ist im Prinzip die ``rechte Alt-Taste'', hat aber auch die Funktion, Zeichen einzugeben, die auf, z.B. einer deutschen Tastatur, sonst nicht zur Verfügung stehen.
    \item[Esc]
        Das ist die sogenannte ``escape''-Taste, auf Deutsch übersetzt ``Flucht''. Mit dieser Taste lassen sich viele Vorgänge abbrechen.
    \item[Cursor]
        Damit sind die Tasten gemeint, mit den man nach rechts, links, oben unten geht. Im Deutschen werden sie auch ``Kursor'', ``Schreibmarke'', ``Eingabezeiger'' usw. bezeichnet, oder einfach als ``Pfeiltasten''.
    \item[Tab]
        Die ``Tabulator''-Taste dient dazu, Text richtig einzurücken.
\end{description}

%%%%%%%%%%%%%%%%%%%%%%%%%%%%%%%%%%%%%%%%%%%%%%%%%%%%%%%%%%%%%%%%%%%%%%%%%%%%%%%%%%%
%%%%%%%%%%%%%%%%%%%%%%%%%%%%%%%%%%%%%%%%%%%%%%%%%%%%%%%%%%%%%%%%%%%%%%%%%%%%%%%%%%%
%%%%%%%%%%%%%%%%%%%%%%%%%%%%%%%%%%%%%%%%%%%%%%%%%%%%%%%%%%%%%%%%%%%%%%%%%%%%%%%%%%%
\subsection{Sonderzeichen}

Die folgenden Sonderzeichen gelten für eine deutsche Tastatur:
\vspace{1em}

\begin{commandtable}{lcl}
\textbf{Tastenkombination} & \textbf{Zeichen} & \textbf{Bedeutung} \\\hline
\endhead

AltGr + Q & @ & At-Zeichen \\
AltGr + 7 & \{ & Geschweifte Klammer auf \\
AltGr + 8 & [ & Eckige Klammer auf \\
AltGr + 9 & ] & Eckige Klammer zu \\
AltGr + 0 & \} & Geschweifte Klammer zu \\
AltGr + ß & \textbackslash & Backslash \\
AltGr + < & | & Pipe \\
AltGr + 3 & \# & Hashtag \\
AltGr + + & \textasciitilde & Tilde \\
% AltGr + 4 & € & Euro-Zeichen \\
% AltGr + E & € & Euro-Zeichen (Alternative) \\
% AltGr + M & µ & Mikro-Zeichen \\
% AltGr + H & ² & Hochgestellt 2 \\
% AltGr + Z & « & Anführungszeichen links (Guillemet) \\
% AltGr + Y & » & Anführungszeichen rechts (Guillemet) \\
\end{commandtable}

%%%%%%%%%%%%%%%%%%%%%%%%%%%%%%%%%%%%%%%%%%%%%%%%%%%%%%%%%%%%%%%%%%%%%%%%%%%%%%%%%%%
%%%%%%%%%%%%%%%%%%%%%%%%%%%%%%%%%%%%%%%%%%%%%%%%%%%%%%%%%%%%%%%%%%%%%%%%%%%%%%%%%%%
%%%%%%%%%%%%%%%%%%%%%%%%%%%%%%%%%%%%%%%%%%%%%%%%%%%%%%%%%%%%%%%%%%%%%%%%%%%%%%%%%%%
\subsection{Referenz}

Die folgende Liste von Tastenkombination gilt für viele, aber nicht für alle Programme:

\begin{commandtable}{lll}
 & \textbf{Tastenkombination} & \textbf{Bedeutung} \\\hline
\endhead
\multicolumn{3}{l}{\textit{Markieren}}\\
& Esc               & Markierung aufheben\\
& Shift+Cursor      & Text markieren (alternativ mit Maus) \\
& Strg+A            & Alles auswählen (``select all'') \\
\\
\multicolumn{3}{l}{\textit{Kopieren und Einfügen}}\\
& Strg+C            & Markierten Text kopieren\\
& Strg+X            & Markierten Text ausschneiden\\
& Strg+V            & Text einfügen\\
& Strg+Z            & Eingabe rückgängig machen (``undo'') \\
\\
\multicolumn{3}{l}{\textit{Häufige Funktionen}}\\
& Strg+S            & Speichern (immer wieder drücken!) \\
& Strg+F            & Suchen (``find'')%\footnote{Ihr werdet vielleicht schon einmal in einem Meme ``Ctrl+F'' gesehen haben. Das bezieht sich in der Tat auf diese Tastenkombination, die allgemein ``suchen'' bedeutet} 
                        \\
& Strg+H            & Suchen und ersetzen (``replace'') \\
\end{commandtable}
